% Test de la relation de conjugaison : 1/OA' en fonction de 1/OA (corrigé de l'exercice 10, chapitre 18)
\begin{tikzpicture}
\begin{axis}[otMesures,otLecture,otAxesBout,width=96mm,height=56mm,xmin=-0.075,xmax=0,ymin=0,ymax=0.11,
  xtick={-0.07,-0.06,...,0},ytick={0,0.02,...,0.1},scaled x ticks=false,scaled y ticks=false,
  xticklabel style={/pgf/number format/fixed,/pgf/number format/precision=2},
  yticklabel style={/pgf/number format/fixed,/pgf/number format/precision=2},
  xlabel={$1/\overline{\mathrm{OA}}$ (\unit{cm^{-1}})},ylabel={$1/\overline{\mathrm{OA'}}$ (\unit{cm^{-1}})}]
\addplot[otNavy,line width=1.1pt,domain=-0.075:0] {x+0.100};
\addplot[only marks,mark=*,mark size=1.6pt,otOrange!85!black] coordinates {(-0.0667,0.0331) (-0.05,0.0505) (-0.04,0.0602) (-0.0333,0.0662) (-0.025,0.0746) (-0.02,0.0806)};
\node[pcLegende,anchor=north west,text=otNavy,align=left] at (axis cs:-0.072,0.105) {droite de pente 1\\ordonnée à l'origine $\dfrac{1}{f'}\approx\qty{0,100}{cm^{-1}}$};
\end{axis}
\end{tikzpicture}
\par\smallskip\legende{Les points s'alignent sur une droite de pente 1 : $\dfrac{1}{\overline{\mathrm{OA'}}}=\dfrac{1}{\overline{\mathrm{OA}}}+\dfrac{1}{f'}$, la relation de conjugaison est vérifiée.}
